\chapter{Metaprogramming and Reflection}\label{ch:metaprogramming}

\index{metaprogramming|(}
\index{reflection|(}

Most of the time, a program manipulates data: numbers, strings, arrays, structs.
But sometimes a program needs to manipulate \emph{itself}---inspect its own structure, examine types at runtime, build code from strings, or interrogate a struct to discover what fields it carries.
These capabilities fall under the umbrella of \textbf{metaprogramming} and \textbf{reflection}.

Lattice is not a ``macro language'' in the tradition of Lisp or Rust.
It does not have compile-time code generation or hygienic macros.
Instead, it offers a pragmatic set of \emph{runtime} introspection tools: functions that let you peek inside values, tokenize source code, convert between structs and maps, and even evaluate strings as Lattice code on the fly.

These tools are powerful, and like any powerful tool, they come with trade-offs.
Let's explore them.

%% ─────────────────────────────────────────────────────────────────────────────
\section{Runtime Code Introspection and Dynamic Features}\label{sec:introspection}
%% ─────────────────────────────────────────────────────────────────────────────

\index{typeof@\lstinline|typeof|}
\index{phase\_of@\lstinline|phase_of|}
\index{to\_string@\lstinline|to_string|}
\index{repr@\lstinline|repr|}

Lattice provides a family of built-in functions for examining values at runtime.
These are not methods on objects---they are global functions available everywhere, defined in \filename{src/eval.c} (for the tree-walk interpreter) and mirrored in the bytecode VM dispatch loops.

\subsection{Type Queries}

The most fundamental introspection is asking ``what \emph{is} this value?''

\begin{lstlisting}[language=Lattice, caption={Querying types at runtime}]
print(typeof(42))           // Int
print(typeof(3.14))         // Float
print(typeof("hello"))      // String
print(typeof(true))         // Bool
print(typeof([1, 2, 3]))    // Array
print(typeof(nil))          // Nil
\end{lstlisting}

The \lstinline|typeof| function returns a string naming the value's type.
Under the hood, it calls \lstinline|builtin_typeof_str()| from \filename{src/builtins.c}, which is a straightforward \lstinline|switch| on the value's internal type tag.
The possible return values are: \lstinline|"Int"|, \lstinline|"Float"|, \lstinline|"Bool"|, \lstinline|"String"|, \lstinline|"Array"|, \lstinline|"Struct"|, \lstinline|"Closure"|, \lstinline|"Unit"|, \lstinline|"Nil"|, \lstinline|"Range"|, \lstinline|"Map"|, \lstinline|"Channel"|, \lstinline|"Buffer"|, \lstinline|"Tuple"|, \lstinline|"Set"|, \lstinline|"Iterator"|, \lstinline|"Ref"|, and \lstinline|"Enum"|.

\begin{defnbox}{typeof()}
\index{typeof@\lstinline|typeof|!definition}
\lstinline|typeof(val)| returns a \lstinline|String| naming the runtime type of \lstinline|val|.
For structs, it returns \lstinline|"Struct"| (not the struct's name---use \lstinline|struct_name()| for that).
\end{defnbox}

This makes \lstinline|typeof| useful for writing generic functions that need to branch on type:

\begin{lstlisting}[language=Lattice, caption={A generic serializer using typeof}]
fn serialize(val: any) -> String {
  let t = typeof(val)
  if t == "Int" || t == "Float" || t == "Bool" {
    return to_string(val)
  }
  if t == "String" {
    return "\"" + val + "\""
  }
  if t == "Array" {
    let parts = val.map(|item| serialize(item))
    return "[" + parts.join(", ") + "]"
  }
  if t == "Nil" {
    return "null"
  }
  return repr(val)
}

print(serialize([1, "hello", true]))
// Output: [1, "hello", true]
\end{lstlisting}

\subsection{Phase Queries}

\index{phase\_of@\lstinline|phase_of|!function}

Lattice's phase system distinguishes fluid (mutable) values from crystallized (frozen) ones.
The \lstinline|phase_of| function lets you check a value's phase at runtime:

\begin{lstlisting}[language=Lattice, caption={Querying value phase}]
flux temperature = 72.0
print(phase_of(temperature))     // fluid

fix pi = 3.14159
print(phase_of(pi))              // crystal

let config = freeze(["debug", "verbose"])
print(phase_of(config))          // crystal
\end{lstlisting}

The return values are \lstinline|"fluid"|, \lstinline|"crystal"|, or \lstinline|"unphased"| (for values that have not been explicitly assigned a phase).
This is particularly useful when writing library code that must respect its callers' mutability constraints:

\begin{lstlisting}[language=Lattice, caption={Respecting caller's phase constraints}]
fn safe_append(collection: any, item: any) -> any {
  if phase_of(collection) == "crystal" {
    // Don't mutate a frozen collection; return a new one
    let thawed = thaw(collection)
    thawed.push(item)
    return freeze(thawed)
  }
  collection.push(item)
  return collection
}
\end{lstlisting}

\subsection{String Representations}

Two functions convert values to strings, each with a different purpose:

\begin{itemize}
  \item \lstinline|to_string(val)| --- produces a human-readable display string.
        Strings are returned without quotes, arrays are formatted with brackets,
        and special values like \lstinline|nil| become \lstinline|"nil"|.
  \item \lstinline|repr(val)| --- produces a representation string that could
        (in principle) be parsed back into the original value. Strings are
        quoted, struct fields are shown, and closures display their parameter
        lists.
\end{itemize}

\begin{lstlisting}[language=Lattice, caption={to\_string vs. repr}]
let greeting = "hello"
print(to_string(greeting))   // hello
print(repr(greeting))        // "hello"

let items = [1, 2, 3]
print(to_string(items))      // [1, 2, 3]
print(repr(items))           // [1, 2, 3]

print(to_string(nil))        // nil
print(repr(nil))             // nil
\end{lstlisting}

\begin{tipbox}{When to Use repr()}
Use \lstinline|repr()| when writing debug output or logging.
It shows you the ``programmer's view'' of a value, including type information that \lstinline|to_string()| strips away.
For user-facing output, prefer \lstinline|to_string()| or string interpolation.
\end{tipbox}

\subsection{Dynamic Evaluation with lat\_eval()}

\index{lat\_eval@\lstinline|lat_eval|}

The most powerful (and most dangerous) introspection tool is \lstinline|lat_eval()|, which takes a string of Lattice source code, parses and evaluates it at runtime, and returns the result:

\begin{lstlisting}[language=Lattice, caption={Dynamic code evaluation}]
let result = lat_eval("2 + 3 * 4")
print(result)  // 14

// Define a function dynamically
lat_eval("fn greet(name: String) { print(\"Hello, \" + name) }")
greet("World")  // Hello, World
\end{lstlisting}

Under the hood, \lstinline|lat_eval()| in \filename{src/eval.c} does exactly what the top-level interpreter does: it tokenizes the string, parses it into an AST, registers any function and struct declarations into the current evaluator's registries, and executes the statements.
Bindings created inside \lstinline|lat_eval()| persist in the caller's scope.

\begin{warnbox}{lat\_eval() and Security}
\lstinline|lat_eval()| executes arbitrary code with full access to the runtime environment.
Never pass untrusted input (user data, network data) to \lstinline|lat_eval()|.
An attacker could read files, make network connections, or execute system commands.
In production code, prefer structured approaches (closures, maps, pattern matching) over \lstinline|lat_eval()|.
\end{warnbox}

\subsection{Completeness Checking with is\_complete()}

\index{is\_complete@\lstinline|is_complete|}

A companion to \lstinline|lat_eval()| is \lstinline|is_complete()|, which checks whether a source string has balanced brackets and parentheses:

\begin{lstlisting}[language=Lattice, caption={Checking source completeness}]
print(is_complete("1 + 2"))        // true
print(is_complete("fn foo() {"))   // false — unclosed brace
print(is_complete("if true { }"))  // true
\end{lstlisting}

This is used internally by the REPL to determine whether to wait for more input or evaluate the current line.
It tokenizes the string, counts bracket depth, and returns \lstinline|true| if the depth is zero or less.

%% ─────────────────────────────────────────────────────────────────────────────
\section{tokenize() --- Inspecting Tokens}\label{sec:tokenize}
%% ─────────────────────────────────────────────────────────────────────────────

\index{tokenize@\lstinline|tokenize|}
\index{lexer!from Lattice code}

The \lstinline|tokenize()| built-in function exposes Lattice's own lexer to user code.
Given a string of source code, it returns an array of \lstinline|Token| structs, each with a \lstinline|type| field (the token type name) and a \lstinline|text| field (the token's text):

\begin{lstlisting}[language=Lattice, caption={Tokenizing Lattice source code}]
let tokens = tokenize("let x = 42 + y")

for tok in tokens {
  print("${tok.type}: ${tok.text}")
}
// Output:
// LET: LET
// IDENT: x
// ASSIGN: ASSIGN
// INT_LIT: 42
// PLUS: PLUS
// IDENT: y
\end{lstlisting}

Each element in the returned array is a struct with two fields:

\begin{itemize}
  \item \lstinline|type| --- a \lstinline|String| naming the token type
        (e.g., \lstinline|"IDENT"|, \lstinline|"INT_LIT"|, \lstinline|"STRING_LIT"|,
        \lstinline|"PLUS"|, \lstinline|"LBRACE"|, \lstinline|"FN"|).
  \item \lstinline|text| --- a \lstinline|String| containing the token's text.
        For identifiers and literals, this is the actual text from the source.
        For operators and keywords, this is the token type name repeated.
\end{itemize}

The trailing \lstinline|EOF| token is stripped from the result, so you only see meaningful tokens.

\subsection{Use Cases for tokenize()}

The most prominent use of \lstinline|tokenize()| is in the self-hosted compiler (see \cref{sec:self-hosted}).
The compiler written in Lattice (\filename{compiler/latc.lat}) calls \lstinline|tokenize()| on the input source file to get its token stream, then implements a recursive-descent parser over those tokens.
This means the self-hosted compiler does not need to reimplement lexing---it reuses the C lexer via this single function call.

Beyond the self-hosted compiler, \lstinline|tokenize()| enables:

\begin{lstlisting}[language=Lattice, caption={Counting keywords in a source file}]
fn count_keywords(source: String) -> Int {
  let tokens = tokenize(source)
  let keywords = ["fn", "let", "flux", "fix", "if", "else",
                   "for", "while", "return", "match", "struct"]
  flux count = 0
  for tok in tokens {
    if tok.type == "IDENT" {
      for kw in keywords {
        if tok.text == kw {
          count = count + 1
        }
      }
    }
  }
  return count
}
\end{lstlisting}

\begin{lstlisting}[language=Lattice, caption={A syntax highlighter skeleton}]
fn highlight(source: String) -> String {
  let tokens = tokenize(source)
  flux result = ""
  for tok in tokens {
    if tok.type == "FN" || tok.type == "LET" || tok.type == "IF" {
      result = result + "\x1b[35m" + tok.text + "\x1b[0m "
    } else if tok.type == "STRING_LIT" {
      result = result + "\x1b[32m\"" + tok.text + "\"\x1b[0m "
    } else if tok.type == "INT_LIT" || tok.type == "FLOAT_LIT" {
      result = result + "\x1b[34m" + tok.text + "\x1b[0m "
    } else {
      result = result + tok.text + " "
    }
  }
  return result
}
\end{lstlisting}

\begin{notebox}{tokenize() and the Bytecode VMs}
The \lstinline|tokenize()| function is available in all three backends.
In the tree-walk interpreter, it calls \lstinline|lexer_tokenize()| directly.
In the bytecode VMs, it is dispatched as a native built-in function that invokes the same C lexer.
The returned tokens are identical regardless of backend.
\end{notebox}

\subsection{Token Types}

Lattice's lexer produces a rich set of token types.
Here is a representative sample:

\begin{table}[h]
\centering
\caption{Common token types returned by tokenize()}\label{tab:token-types}
\begin{tabular}{@{} l l p{5.5cm} @{}}
\toprule
\textbf{Token Type} & \textbf{Example Text} & \textbf{Description} \\
\midrule
\texttt{IDENT}       & \texttt{my\_var}    & Identifier \\
\texttt{INT\_LIT}    & \texttt{42}         & Integer literal \\
\texttt{FLOAT\_LIT}  & \texttt{3.14}       & Float literal \\
\texttt{STRING\_LIT} & \texttt{hello}      & String literal (without quotes) \\
\texttt{FN}          & \texttt{FN}         & \lstinline|fn| keyword \\
\texttt{LET}         & \texttt{LET}        & \lstinline|let| keyword \\
\texttt{FLUX}        & \texttt{FLUX}       & \lstinline|flux| keyword \\
\texttt{FIX}         & \texttt{FIX}        & \lstinline|fix| keyword \\
\texttt{IF}          & \texttt{IF}         & \lstinline|if| keyword \\
\texttt{PLUS}        & \texttt{PLUS}       & \texttt{+} operator \\
\texttt{ASSIGN}      & \texttt{ASSIGN}     & \texttt{=} assignment \\
\texttt{LBRACE}      & \texttt{LBRACE}     & \texttt{\{} opening brace \\
\texttt{RBRACE}      & \texttt{RBRACE}     & \texttt{\}} closing brace \\
\bottomrule
\end{tabular}
\end{table}

For a complete list, see \cref{appA-grammar-reference}.

%% ─────────────────────────────────────────────────────────────────────────────
\section{Struct Reflection Functions}\label{sec:struct-reflection}
%% ─────────────────────────────────────────────────────────────────────────────

\index{struct reflection|(}
\index{struct\_name@\lstinline|struct_name|}
\index{struct\_fields@\lstinline|struct_fields|}
\index{struct\_to\_map@\lstinline|struct_to_map|}
\index{struct\_from\_map@\lstinline|struct_from_map|}

Lattice structs are not opaque---at runtime, you can interrogate them to discover their name, fields, and values.
This makes it possible to write generic code that operates on \emph{any} struct, regardless of its definition.

\subsection{struct\_name()}

\begin{lstlisting}[language=Lattice, caption={Getting a struct's type name}]
struct User {
  name: String,
  age: Int,
}

let alice = User { name: "Alice", age: 30 }
print(struct_name(alice))  // User
\end{lstlisting}

The \lstinline|struct_name()| function returns the struct's type name as a string.
This is distinct from \lstinline|typeof(alice)|, which returns \lstinline|"Struct"| for all struct instances.
With \lstinline|struct_name()| you can distinguish between different struct types at runtime.

\subsection{struct\_fields()}

\begin{lstlisting}[language=Lattice, caption={Listing a struct's field names}]
struct Point {
  x: Float,
  y: Float,
}

let origin = Point { x: 0.0, y: 0.0 }
let fields = struct_fields(origin)
print(fields)  // [x, y]
\end{lstlisting}

The \lstinline|struct_fields()| function returns an array of field name strings.
The names are in declaration order---the same order you specified when defining the struct.

\subsection{struct\_to\_map()}

\index{struct\_to\_map@\lstinline|struct_to_map|!function}

Perhaps the most useful reflection function, \lstinline|struct_to_map()| converts a struct instance into a \lstinline|Map| where the keys are field names and the values are deep clones of the field values:

\begin{lstlisting}[language=Lattice, caption={Converting a struct to a map}]
struct Config {
  host: String,
  port: Int,
  debug: Bool,
}

let cfg = Config { host: "localhost", port: 8080, debug: true }
let map = struct_to_map(cfg)

print(map["host"])    // localhost
print(map["port"])    // 8080
print(map["debug"])   // true
\end{lstlisting}

This is invaluable for serialization.
Instead of writing a custom serializer for each struct type, you can convert to a map and pass it to \lstinline|json_stringify()|:

\begin{lstlisting}[language=Lattice, caption={Generic struct serialization via struct\_to\_map}]
fn to_json(val: any) -> String {
  if typeof(val) == "Struct" {
    let map = struct_to_map(val)
    return json_stringify(map)
  }
  return json_stringify(val)
}

struct Product {
  name: String,
  price: Float,
  in_stock: Bool,
}

let item = Product { name: "Widget", price: 9.99, in_stock: true }
print(to_json(item))
// Output: {"name":"Widget","price":9.99,"in_stock":true}
\end{lstlisting}

\subsection{struct\_from\_map()}

\index{struct\_from\_map@\lstinline|struct_from_map|!function}

The inverse of \lstinline|struct_to_map()| is \lstinline|struct_from_map()|, which creates a struct instance from a type name and a map of field values:

\begin{lstlisting}[language=Lattice, caption={Creating a struct from a map}]
struct User {
  name: String,
  age: Int,
}

let data = Map::new()
data["name"] = "Bob"
data["age"] = 25

let user = struct_from_map("User", data)
print(user.name)  // Bob
print(user.age)   // 25
\end{lstlisting}

The struct type must already be defined in the current scope.
If a field is missing from the map, it defaults to \lstinline|nil|.
This function is particularly useful for deserialization---parsing JSON into a map, then converting the map into a typed struct:

\begin{lstlisting}[language=Lattice, caption={JSON deserialization into a struct}]
struct Task {
  title: String,
  done: Bool,
  priority: Int,
}

fn parse_task(json_str: String) -> any {
  let map = json_parse(json_str)
  return struct_from_map("Task", map)
}

let task = parse_task("{\"title\": \"Write docs\", \"done\": false, \"priority\": 1}")
print(task.title)     // Write docs
print(task.priority)  // 1
\end{lstlisting}

\subsection{Combining Reflection for Generic Programming}

With these four functions, you can write powerful generic code:

\begin{lstlisting}[language=Lattice, caption={A generic equality checker for structs}]
fn struct_equal(a: any, b: any) -> Bool {
  if typeof(a) != "Struct" || typeof(b) != "Struct" {
    return false
  }
  if struct_name(a) != struct_name(b) {
    return false
  }
  let fields_a = struct_fields(a)
  let fields_b = struct_fields(b)
  if len(fields_a) != len(fields_b) {
    return false
  }
  let map_a = struct_to_map(a)
  let map_b = struct_to_map(b)
  for field in fields_a {
    if map_a[field] != map_b[field] {
      return false
    }
  }
  return true
}

struct Point {
  x: Float,
  y: Float,
}

let p1 = Point { x: 1.0, y: 2.0 }
let p2 = Point { x: 1.0, y: 2.0 }
let p3 = Point { x: 3.0, y: 4.0 }

print(struct_equal(p1, p2))  // true
print(struct_equal(p1, p3))  // false
\end{lstlisting}

\begin{lstlisting}[language=Lattice, caption={A generic struct printer}]
fn describe_struct(val: any) -> String {
  if typeof(val) != "Struct" {
    return repr(val)
  }
  let name = struct_name(val)
  let fields = struct_fields(val)
  let map = struct_to_map(val)
  flux parts = []
  for field in fields {
    parts.push(field + ": " + repr(map[field]))
  }
  return name + " { " + parts.join(", ") + " }"
}

struct Color {
  r: Int,
  g: Int,
  b: Int,
}

let coral = Color { r: 255, g: 127, b: 80 }
print(describe_struct(coral))
// Output: Color { r: 255, g: 127, b: 80 }
\end{lstlisting}

\index{struct reflection|)}

%% ─────────────────────────────────────────────────────────────────────────────
\section{format() --- String Building}\label{sec:format}
%% ─────────────────────────────────────────────────────────────────────────────

\index{format@\lstinline|format|}
\index{string formatting}

While string interpolation with \lstinline|${}| handles most formatting needs, the \lstinline|format()| function provides a positional placeholder approach similar to Python's \lstinline|str.format()| or Rust's \lstinline|format!()| macro.

\subsection{Basic Usage}

\begin{lstlisting}[language=Lattice, caption={Using format() for string building}]
let msg = format("{} is {} years old", "Alice", 30)
print(msg)  // Alice is 30 years old

let report = format("Total: {} items, {} in stock", 150, 42)
print(report)  // Total: 150 items, 42 in stock
\end{lstlisting}

The format string uses \lstinline|{}| as a placeholder.
Arguments are consumed left to right, one per placeholder.
Each argument is converted to its display string (equivalent to calling \lstinline|to_string()| on it).

\subsection{Escaping Braces}

If you need a literal \lstinline|{| or \lstinline|}| in the output, double it:

\begin{lstlisting}[language=Lattice, caption={Escaping braces in format strings}]
let json_template = format("{{\"name\": \"{}\", \"age\": {}}}", "Bob", 25)
print(json_template)  // {"name": "Bob", "age": 25}
\end{lstlisting}

The implementation in \filename{src/format\_ops.c} handles this by checking for consecutive \lstinline|{{| and \lstinline|}}| patterns, emitting a single brace for each pair.

\subsection{Error Handling}

If you provide fewer arguments than placeholders, \lstinline|format()| raises a runtime error:

\begin{lstlisting}[language=Lattice, caption={format() error on missing arguments}]
// This will raise an error:
// let msg = format("{} and {}", "Alice")
// Error: format: not enough arguments for placeholders
\end{lstlisting}

Extra arguments are silently ignored---only the placeholders that exist in the format string consume arguments.

\subsection{format() vs. String Interpolation}

When should you use \lstinline|format()| instead of string interpolation?

\begin{itemize}
  \item \textbf{Dynamic format strings}: When the template is not known at compile time---for instance, loaded from a configuration file or received over a network.
  \item \textbf{Reusable templates}: When you want to define a template once and apply it with different arguments.
  \item \textbf{Array-based arguments}: When the arguments come from a computed source rather than inline variables.
\end{itemize}

\begin{lstlisting}[language=Lattice, caption={Dynamic templates with format()}]
// A logging system with configurable format
flux log_template = "[{}] {}: {}"

fn log(level: String, message: String) {
  let timestamp = to_string(time_ms())
  print(format(log_template, timestamp, level, message))
}

log("INFO", "Server started")
// Output: [1709234567890] INFO: Server started

log("ERROR", "Connection refused")
// Output: [1709234567891] ERROR: Connection refused
\end{lstlisting}

\begin{tipbox}{Prefer Interpolation for Readability}
For most cases, string interpolation (\lstinline|"Hello, ${name}!"|) is more readable than \lstinline|format("Hello, {}!", name)|.
Reserve \lstinline|format()| for situations where the template itself is dynamic or where you are building strings programmatically.
\end{tipbox}

%% ─────────────────────────────────────────────────────────────────────────────
\section{The Power and Danger of Runtime Features}\label{sec:power-danger}
%% ─────────────────────────────────────────────────────────────────────────────

\index{lat\_eval@\lstinline|lat_eval|!dangers}
\index{metaprogramming!best practices}

Runtime introspection and dynamic evaluation are a double-edged sword.
Let's be honest about both edges.

\subsection{The Power}

Dynamic features let you build things that would otherwise require code generation or preprocessors:

\begin{lstlisting}[language=Lattice, caption={A configuration-driven plugin system}]
struct Plugin {
  name: String,
  version: String,
  init: fn() -> Bool,
}

fn load_plugin_from_config(config: Map) -> any {
  let name = config["name"]
  let source = read_file("plugins/" + name + ".lat")
  lat_eval(source)
  // The plugin registers itself via a global function
  return true
}
\end{lstlisting}

\begin{lstlisting}[language=Lattice, caption={A test runner that discovers tests}]
fn find_test_functions(source: String) -> Array {
  let tokens = tokenize(source)
  flux test_names = []
  for i in 0..len(tokens) - 1 {
    if tokens[i].type == "TEST" {
      if i + 1 < len(tokens) && tokens[i + 1].type == "STRING_LIT" {
        test_names.push(tokens[i + 1].text)
      }
    }
  }
  return test_names
}

let source = read_file("my_tests.lat")
let tests = find_test_functions(source)
print("Found " + to_string(len(tests)) + " tests:")
for name in tests {
  print("  - " + name)
}
\end{lstlisting}

\begin{lstlisting}[language=Lattice, caption={A generic ORM-style mapper}]
fn map_rows_to_structs(rows: Array, struct_type: String) -> Array {
  flux results = []
  for row in rows {
    let instance = struct_from_map(struct_type, row)
    results.push(instance)
  }
  return results
}

struct Employee {
  name: String,
  department: String,
  salary: Float,
}

// Imagine these rows came from a database query
let rows = [
  Map::new(),
  Map::new(),
]
rows[0]["name"] = "Alice"
rows[0]["department"] = "Engineering"
rows[0]["salary"] = 95000.0
rows[1]["name"] = "Bob"
rows[1]["department"] = "Marketing"
rows[1]["salary"] = 85000.0

let employees = map_rows_to_structs(rows, "Employee")
for emp in employees {
  print("${emp.name} in ${emp.department}")
}
// Output:
// Alice in Engineering
// Bob in Marketing
\end{lstlisting}

\subsection{The Dangers}

Every runtime feature trades away some combination of safety, performance, and readability.

\begin{description}
  \item[Performance.] \lstinline|lat_eval()| must lex, parse, and interpret the string on every call.
    There is no caching and no compilation---it is a full interpreter invocation.
    In a tight loop, this overhead is devastating.

  \item[Security.] \lstinline|lat_eval()| on untrusted input is a code injection vulnerability.
    The evaluated code runs with the same privileges as the host program.

  \item[Debuggability.] Dynamically generated code does not have source file locations.
    Error messages will reference the evaluated string, not a file and line number.
    Stack traces from \lstinline|lat_eval()| are harder to read.

  \item[Type safety.] Reflection functions like \lstinline|struct_to_map()| erase type information.
    The resulting map has string keys and \lstinline|any|-typed values---you lose the compiler's ability
    to check field names and types.

  \item[Refactoring.] Code that uses \lstinline|lat_eval()| to define functions or structs dynamically
    is invisible to static analysis tools, formatters, and documentation generators.
    Renaming a function in your editor will not catch references inside evaluated strings.
\end{description}

\begin{warnbox}{The lat\_eval() Rule of Thumb}
If you find yourself reaching for \lstinline|lat_eval()|, stop and ask: ``Can I solve this with closures, maps, or pattern matching instead?''
The answer is almost always yes.
Reserve \lstinline|lat_eval()| for genuinely dynamic scenarios: plugin systems, REPLs, code generators, and development tools.
\end{warnbox}

\subsection{Guidelines for Responsible Metaprogramming}

\begin{enumerate}
  \item \textbf{Prefer closures over eval.}
        Instead of building a function as a string and evaluating it,
        accept a closure parameter and call it.

  \item \textbf{Use struct reflection for serialization, not for business logic.}
        \lstinline|struct_to_map()| and \lstinline|struct_from_map()| are ideal
        for converting between data formats. They are less ideal for implementing
        core application behavior.

  \item \textbf{Validate before evaluating.}
        If you must use \lstinline|lat_eval()|, validate the input first.
        At minimum, check that it tokenizes without errors using
        \lstinline|is_complete()|.

  \item \textbf{Use tokenize() for analysis, not execution.}
        \lstinline|tokenize()| is safe---it only reads, never evaluates.
        It is the right tool for syntax highlighting, code analysis, and
        documentation generation.

  \item \textbf{Document dynamic behavior.}
        When your code uses reflection or dynamic evaluation, add comments
        explaining \emph{why} the dynamic approach is necessary and what
        the expected inputs are.
\end{enumerate}

\begin{lstlisting}[language=Lattice, caption={Closures instead of eval: the better way}]
// BAD: building code as strings
// let code = "fn process(x: Int) -> Int { return x * " + to_string(multiplier) + " }"
// lat_eval(code)

// GOOD: use a closure that captures the value
fn make_processor(multiplier: Int) -> fn(Int) -> Int {
  return |x| x * multiplier
}

let process = make_processor(3)
print(process(10))  // 30
\end{lstlisting}

%% ─────────────────────────────────────────────────────────────────────────────
\section{Exercises}\label{sec:meta-exercises}
%% ─────────────────────────────────────────────────────────────────────────────

\begin{enumerate}
  \item Write a function \lstinline|pretty_print(val: any)| that produces a formatted, indented string representation of any value.
        For structs, use \lstinline|struct_name()| and \lstinline|struct_to_map()| to show the struct type and fields.
        For arrays and maps, indent nested elements.

  \item Use \lstinline|tokenize()| to write a function that counts the number of functions defined in a Lattice source file.
        Look for the \lstinline|FN| token followed by an \lstinline|IDENT| token.

  \item Write a \lstinline|diff_structs(a: any, b: any)| function that compares two struct instances of the same type and returns an array of field names that differ.
        Use \lstinline|struct_fields()| and \lstinline|struct_to_map()| to implement it generically.

  \item Create a simple calculator REPL that reads a line from the user with \lstinline|input()|, evaluates it with \lstinline|lat_eval()|, and prints the result.
        Add basic error handling: catch evaluation errors and display a helpful message instead of crashing.

  \item (Challenge) Write a ``struct validator'' function that takes a struct instance and a map of \lstinline|{field_name: validation_closure}| pairs.
        For each field, call the corresponding validation closure with the field's value and collect any failures.
        Return an array of error messages for fields that failed validation.
\end{enumerate}

%% ─────────────────────────────────────────────────────────────────────────────
\vspace{1em}
\noindent\textbf{What's Next}\quad
We have peered behind the curtain at Lattice's runtime introspection capabilities---type queries, tokenization, struct reflection, and dynamic evaluation.
With the ``Under the Hood'' part complete, we are ready to put everything together.
In \cref{ch:building-cli-tool}, we will build a complete command-line tool from scratch, applying the language features, concurrency patterns, and standard library knowledge we have accumulated throughout this book.

\index{metaprogramming|)}
\index{reflection|)}
